% Sección: Estrategia de modelado
%
% Mencionar, no sé en qué orden: partición en train y test; preprocesamiento de
% los datos, normalización, pca, etc.; existencia de hiperparámetros y el método
% utilizado para seleccionarlos, CV con 5 folds; ¿capaz aquí sí el software,
% paquetes de R, etc.?; procedimiento general con cada modelo.

Como se mencionó en la sección anterior, los datos originales se dividieron en
un conjunto de entrenamiento, con aproximadamente un 70\% de las observaciones,
y uno de testeo o evaluación con el resto. El primero se utilizó para entrenar
y validar los modelos, y el segundo para calcular las métricas en base a las
que se eligieron los dos mejores.

En términos generales, la estrategia de modelado debió tomar en cuenta los
siguientes problemas. En primer lugar, el análisis exploratorio mostró que los
variables predictoras tienen unidades dispares y pueden llegar a presentar
correlaciones lineales muy altas. Esto llevó a la decisión de normalizar las
variables y de probar estrategias de reducción de la dimensión o regularización
en algunos casos. En segundo lugar, varios de los modelos propuestos
(enumerados en la Subsección \ref{sec:modelos}) cuentan con uno o más
hiperparámetros cuyos valores deben ser establecidos antes de la evaluación en
los datos de testeo. En algunos casos se eligieron determinados valores en
función de lo observado en el proceso de modelado y del costo computacional. En
los casos en los que no, dichos valores se eligieron por validación cruzada de
cinco folds buscando sobre grillas de valores definidas de antemano. Esto
introdujo el problema de seleccionar una métrica en base a la cual elegir los
valores más adecuados de los hiperparámetros. Por motivos que se detallan en la
Subsección \ref{sec:metricas} se eligió una métrica distinta a la de evaluación
de los modelos. Por último, se debió establecer un umbral de predicción que
determine qué casos se clasifican como malignos y cuáles como benignos.

El trabajo se realizó íntegramente en R 4.5.2 y RStudio 2025.09.2
\parencite{rcoreteam_2025_language,positteam_2025_rstudio}, fundamentalmente
con los conjuntos de paquetes de \verb|{tidyverse}| y \verb|{tidymodels}|
\parencite{wickham_etal_2019_welcome,kuhn_wickham_2020_tidymodels}. Estos
ofrecen una interfaz sumamente consistente para gran parte del proceso de
modelado, lo cual simplificó mucho el lidiar con seis modelos distintos y todos
los pasos que implica su construcción. En cuanto a los modelos en sí,
\verb|{tidymodels}| hace uso de otros paquetes para implementarlos. En este
caso se usaron los paquetes \verb|{ranger}|, \verb|{glmnet}|, \verb|{kknn}| y
\verb|{discrim}|
\parencite{wright_ziegler_2017_ranger,friedman_etal_2010_regularization,tay_etal_2023_elastic,schliep_hechenbichler_2025_kknn}.

\subsection{Métricas de evaluación y validación}\label{sec:metricas}

% Aclarar que hay dos métricas usadas, área bajo la curva Precision-Recall para
% elegir hiperparámetros, y puntaje $F_\beta$ para evaluar el desempeño en los
% datos de testeo. ¿Otras métricas como precision y recall? Justificar $F_\beta$
% con riesgo de no detectar cáncer vs. riesgo de diagnosticar erróneamente a
% alguien. Mencionar problema con parámetros $\beta$ y el umbral de decisión.
% Explicar cómo se elige el umbral en los modelos. ¿$\beta = 1$?¿Cómo justificar
% otra ponderación?

La elección de los dos mejores modelos se hizo en base al puntaje $F_\beta$
(definido en los párrafos siguientes) con $\beta = 2$ y calculado sobre el
conjunto de testeo. Sin embargo, también se calculó la \textit{accuracy}, la
precisión, el \textit{recall} y el área bajo la curva \textit{precision-recall}
para poder interpretar mejor los resultados. La selección del puntaje $F_2$
como métrica principal se basó en la consideración de las consecuencias de los
falsos positivos y los falsos negativos en la identificación de tumores
malignos. Durante todo el ejercicio se consideró como clase positiva a los
tumores malignos.

\begin{table}[h]
\begin{tabular*}{\textwidth}[t]{@{\extracolsep{\fill}}lcc}
\toprule
\multirow{2}{3em}{Realidad}  & \multicolumn{2}{c}{Predicción} \\
            \cmidrule{2-3}
          & Maligno                 & Benigno \\
\midrule
Maligno   & Verdadero Positivo (VP) & Falso Negativo (FN) \\
Benigno   & Falso Positivo (FP)     & Verdadero Negativo (VN) \\
\bottomrule
\end{tabular*}
\caption{Matriz de confusión.}
\label{tab:confusion}
\end{table}

El Cuadro \ref{tab:confusion} muestra la matriz de confusión que corresponde al
problema en cuestión. Un falso negativo, en este escenario, es el peor de los
resultados ya que implica que un tumor maligno no fue detectado. Si bien el
diagnóstico definitivo del cáncer no se obtiene mediante un único método, su
detección temprana es fundamental para aumentar las chances de supervivencia de
las pacientes y cualquier retraso pone en riesgo sus vidas
\parencite{saunders_jassal_2009_breast}. Un falso positivo, si bien no alcanza
el mismo nivel de gravedad, sigue sin ser un buen resultado; significa que se
identificó como maligno un tumor benigno, lo cual puede llevar a que la persona
deba someterse a estudios invasivos como biopsias o a otros tratamientos. Como
estamos lidiando con tumores mamarios, esto puede significar cirugías,
radioterapia o quimioterapia \parencite{saunders_jassal_2009_breast},
dependiendo de la gravedad del error. Como es razonable en cualquier problema
de clasificación, el objetivo fue minimizar los casos mal clasificados, salvo
que, en este caso, hay un tipo de error que es prioritario evitar por sobre el
otro.

Dos métricas habituales que sirven en este caso son la precisión y el
\textit{recall}. El \textit{recall} o \textit{true positve rate} mide la
proporción de los casos verdaderamente positivos que fueron predichos como
positivos por el modelo:

\begin{equation}
\text{Recall} = \frac{\text{VP}}{\text{VP} + \text{FN}}
\end{equation}

Por su parte, la precisión o \textit{positive predictive value} mide la
proporción de los predichos como positivos por el modelo que realmente son
positivos:

\begin{equation}
\text{Precision} = \frac{\text{VP}}{\text{VP} + \text{FP}}
\end{equation}

Ambas métricas se aproximan a 0 a medida que aumentan los casos mal
clasificados, y a 1 a medida que disminuyen. Sin embargo, existe cierto
\textit{trade-off} entre ambas; si predecimos todos los casos como positivos,
entonces no habrá falsos negativos y el \textit{recall} será 1, pero
inevitablemente habrá falsos positivos, por lo cual la precisión se acercará a
0. Además, puede suceder que la proporción de casos realmente positivos sobre
el total sea muy alto, y que, por lo tanto, un mal modelo logre una alta
precisión simplemente porque hay pocos casos realmente negativos para ser
clasificados como positivos. En ese caso seguramente haya más falsos negativos
y el \textit{recall} será alto.

El puntaje $F_\beta$ ofrece un balance entre ambas métricas. Es su media
armónica ponderada de forma tal que el \textit{recall} pesa $\beta$ veces más
que la precisión:

\begin{equation}
\begin{aligned}
  F_\beta &= \frac{(1 + \beta^2) \text{Precision} \times  \text{Recall}}{\beta^2 \text{Precision} + \text{Recall}} \\
  &= \frac{(1 + \beta^2) \text{VP}}{(1 + \beta^2) \text{VP} + \beta^2 \text{FN} + \text{FP}}
\end{aligned}
\end{equation}

La elección del valor de $\beta$ representó un problema difícil de resolver en
el trabajo. En principio parece evidente que un valor de $\beta = 1$ es
inadecuado ya que las consecuencias de los falsos negativos son más graves que
las de los falsos positivos. Sin embargo, no es claro cómo ir más allá de un
orden en la gravedad de los errores. Un antecedente del campo de la
bioinformática que utiliza un valor de $\beta = 2$ es
\textcite{sanchez-herrero_etal_2024_forecasting}. Sin embargo, otros antecedentes
indican que lo más común es usar $\beta = 1$, y que, de todas formas, el
puntaje $F_\beta$ es inadecuado en casos clínicos
\parencite{assel_vickers_2025_score}. Finalmente se optó por usar un valor de
$\beta = 2$, que es relativamente arbitrario pero mayor a 1.

Las tres medidas presentadas requieren que se establezca un umbral de
predicción que permita convertir las probabilidades de pertenecia a la clase
maligna que estiman los modelos en una clasificación de las observaciones en
dos clases. Este umbral se seleccionó durante la etapa de validación de los
modelos en los datos de entrenamiento (esto se desarrollará en más detalle en
la Sección \ref{sec:detalle}).

Como se mencionó anteriormente, para elegir los hiperparámetros de algunos
modelos se realizó una validación cruzada sobre cinco folds. La métrica que se
buscó maximizar en este caso fue el área bajo la curva
\textit{precision-recall}, no el puntaje $F_2$, ya que la primera no depende
del umbral de decisión y también refleja un balance entre ambas medidas. De lo
contrario, se debería haber incluido el umbral como hiperparámetro a optimizar
durante la validación cruzada, lo cual hubiese complejizado la búsqueda.

\subsection{Modelos a implementar}\label{sec:modelos}

% Enumerar y describir brevemente los modelos a implementar. Mencionar los
% hiperparámetros que tiene cada uno, cuáles van a ser optimizados y cuáles no.
% Justificar los valores utilizados para los que no se optimizan.

Los modelos que se implementaron fueron los siguientes:

\begin{enumerate}

  \item Regresión logística sobre las componentes principales de las variables
    predictoras normalizadas. Se optimizó el número de componentes a utilizar
    mediante validación cruzada. Se decidió hacer una reducción de la dimensión
    dada la alta correlación que se encontró entre las variables predictoras
    durante el análisis exploratorio.
  \item Regresión logística con penalización LASSO sobre las variables
    predictoras originales normalizadas. Se optimizó el hiperparámetro de
    penalización $\lambda$ mediante validación cruzada. Se incluyó este modelo
    ya que la penalización LASSO no sólo permite seleccionar variables
    relevantes, sino que también enfrenta el problema de la correlación entre
    las predictoras limitando el tamaño de los coeficientes.
  \item Regresión logística con penalización LASSO sobre las componentes
    principales de las variables predictoras normalizadas. Al igual que en el
    modelo anterior, se optimizó el hiperparámetro $\lambda$. Se incluyó este
    modelo como método alternativo de selección de componentes principales; el
    primer modelo selecciona las primeras $n$ componentes, mientras que este
    selecciona $m$ componentes que no necesariamente son las primeras.
  \item Bosques aleatorios sobre las variables predictoras originales
    normalizadas. Este modelo cuenta con tres hiperparámetros (al menos en la
    implementación usada): la cantidad de árboles por bosque, la cantidad
    mínima de observaciones por hoja, y la cantidad de variables a utilizar por
    \textit{split} en cada árbol. Durante la validación se notó que la cantidad
    de árboles no afectaba notoriamente el desempeño del modelo, por lo cual se
    eligió usar 100 árboles y se optimizaron los otros dos parámetros (sobre
    grillas reducidas para no aumentar demasiado el costo computacional).
  \item Análisis discriminante lineal sobre las componentes principales de las
    variables predictoras normalizadas. Se optimizó el número de componentes a
    incluir en el modelo. Se utilizaron las componentes principales ya que
    durante el modelado se observó que usando sólo las variables normalizadas
    el modelo tenía un desempeño mucho peor que el del resto.
  \item $K$ vecinos más cercanos sobre las variables predictoras originales
    normalizadas. Se optimizó la cantidad de vecinos.

\end{enumerate}
